\documentclass[12pt]{article}
\usepackage[a4paper,textwidth=165.1mm,top=19mm,bottom=17mm]{geometry}
\usepackage{newtxtext,amsmath,amssymb,newtxmath,graphicx,booktabs,xcolor,fancyhdr,float}
\usepackage[hidelinks]{hyperref}
\setlength{\headheight}{15pt}
\pagestyle{fancy}\fancyhf{}\lhead{Synthetic workflow rehearsal}\rhead{\thepage}
\setlength{\parindent}{0pt}\setlength{\parskip}{4pt}
\setlength{\intextsep}{6pt}
\setlength{\abovecaptionskip}{6pt}
\raggedbottom
\input{numbers.tex}
\begin{document}
\begin{center}
{\LARGE\bfseries Growth, Loss and Numerical Evidence\par}
\smallskip
{\large A Controlled Modeling Rehearsal\par}
{\normalsize Codex MCM Adapter Demonstration\par}
\end{center}
{\small
\textbf{Abstract.}
We compare fourth-order Runge--Kutta (RK4) with a closed-form reference for a
synthetic biomass model. At time 30, baseline biomass is \BaselineBiomass\ units.
The largest sampled baseline error is \mbox{$\SolverError$} units. Nine prescribed
growth and loss combinations produce terminal biomass from \MinimumBiomass\ to
\MaximumBiomass\ units. The resulting envelope describes the range over those
settings. The rehearsal uses synthetic inputs throughout.\par}

\section{Problem Restatement and Analysis}
\label{sec:problem}
The rehearsal asks whether RK4 matches a known trajectory, how growth and loss
affect terminal biomass, and how far the trajectories diverge.
Figure~\ref{fig:workflow} traces the inputs and verification steps.
\begin{figure}[H]\centering
\includegraphics{figures/workflow.pdf}
\caption{Evidence flow for the synthetic rehearsal. RK4 and the analytic reference
share the model contract. CSVs supply plots and paper numbers.
The data manifest records their SHA-256 hashes.}\label{fig:workflow}
\end{figure}

\section{Assumptions and Notation}
\label{sec:assumptions}
The system is spatially homogeneous with constant coefficients and nonnegative
biomass. Table~\ref{tab:assumptions} defines the baseline and the units used in
the synthetic calculations.
\begin{table}[H]\centering
\caption{Baseline quantities and units from the synthetic model contract.}
\label{tab:assumptions}
\begin{tabular}{lll}\toprule
Symbol & Meaning & Baseline value\\\midrule
$B_0$ & Initial biomass & $12$ biomass units\\
$r$ & Intrinsic growth rate & $0.35$ per time unit\\
$d$ & Constant per-capita loss rate & $0.08$ per time unit\\
$K$ & Growth carrying-capacity parameter & $100$ biomass units\\
$T$ & Time horizon & $30$ time units\\
$h$ & Integration step & $0.1$ time units\\
\bottomrule\end{tabular}
\end{table}

\clearpage
\section{Model Construction and Solution}
\label{sec:model}
The balance between density-limited growth and loss gives
\begin{equation}\dot B=rB\left(1-\frac{B}{K}\right)-dB
             =aB-qB^2,\qquad a=r-d,\quad q=r/K.\end{equation}
For $a\ne0$ and $B_0>0$, substituting $u=1/B$ gives $\dot u+au=q$ and hence
\begin{equation} B(t)=\frac{B_0e^{at}}{1+(qB_0/a)(e^{at}-1)}.\end{equation}
At $a=0$, integration yields $B(t)=B_0/(1+qB_0t)$. With $B_0=0$, biomass stays
zero. RK4 uses a shortened final step when the horizon falls between grid points.

The comparison uses $E_h=\max_i|B_h(t_i)-B_{\rm exact}(t_i)|$ on the solver grid.
At $h=0.1$, the error is $\SolverError$ biomass units, below the contract's
$10^{-3}$ threshold. Halving $h$ from $0.2$ to $0.1$ and again to $0.05$ reduces
the error by factors of \FirstErrorRatio\ and \SecondErrorRatio, close to the
fourth-order factor of 16. Figure~\ref{fig:trajectory} shows where the numerical
and analytic solutions differ over time.
\begin{figure}[H]\centering
\includegraphics{figures/trajectory.pdf}
\caption{Baseline numerical verification at all 301 points with $h=0.1$.
(a) The RK4 and analytic biomass trajectories overlap at the displayed scale.
(b) The signed difference is shown in units of $10^{-8}$ biomass.}
\label{fig:trajectory}
\end{figure}

\clearpage
\section{Sensitivity Analysis}
\label{sec:sensitivity}
We vary $r\in\{0.28,0.35,0.42\}$ and $d\in\{0.04,0.08,0.12\}$, keeping $K$, $B_0$
and the horizon fixed. Figure~\ref{fig:sensitivity} shows that terminal biomass
increases with growth for each loss rate and decreases with loss for each growth
rate. Across the nine settings, it ranges from \MinimumBiomass\ to
\MaximumBiomass\ units.

For $r>d$, the positive equilibrium and its local derivatives are
\begin{align}
B^*&=K\left(1-\frac{d}{r}\right),\\
\frac{\partial B^*}{\partial d}&=-\frac{K}{r},
&\frac{\partial B^*}{\partial r}&=\frac{Kd}{r^2}.
\end{align}
These equilibrium derivatives agree with the directions observed at the finite
horizon. The finite-time values depend on the initial state and the time
available to approach equilibrium.
\begin{figure}[H]\centering
\includegraphics{figures/sensitivity.pdf}
\caption{Terminal biomass at $T=30$ across the nine tested growth and loss
combinations. Each cell contains one RK4 result rounded to two decimals.
The outlined center cell is the baseline.}\label{fig:sensitivity}
\end{figure}

\clearpage
\section{Model Evaluation and Recommendations}
\label{sec:evaluation}
The analytic comparison verifies RK4 accuracy. The parameter envelope uses the
pointwise extrema of the nine trajectories in the prescribed set $\mathcal S$.
\begin{equation}
[L(t),U(t)]=\left[\min_{s\in\mathcal S}B_s(t),\;
                      \max_{s\in\mathcal S}B_s(t)\right],\quad |\mathcal S|=9.
\end{equation}
Figure~\ref{fig:envelope} shows the spread developing from a shared initial value.
The bounds apply to the nine runs. Exceedance probabilities would require a
distribution over inputs or empirical observations.
\begin{figure}[H]\centering
\includegraphics{figures/envelope.pdf}
\caption{Pointwise minimum and maximum of nine trajectories on the common 301-point grid.
Dashed lines show individual scenarios and the solid line marks the baseline.}
\label{fig:envelope}
\end{figure}

Before interpreting the scenario results, verify numerical accuracy and report the
parameter grid. An empirical application would also need calibration and an
assessment of spatial effects.

\section*{Report on Use of AI Tools}
Codex assisted with the model contract, code, figures, manuscript and verification.

\appendix\section{Reproduction}
\label{sec:reproduction}
The example README provides the test and TeX build commands. Frozen CSVs, file
hashes, and figure provenance accompany the manuscript in the evidence directory.
\end{document}
